% =====================================================================
% บทที่ 6 การดูแลระบบ (ผู้อ่าน: ผู้ดูแลระบบ)
% =====================================================================
\label{ch:admin}

\begin{guide}
บทนี้สำหรับผู้ที่ดูแลระบบหลังติดตั้ง เช่น เจ้าหน้าที่ไอทีของหน่วยงานที่รับระบบไปใช้
\begin{itemize}
  \item ตารางงานประจำและความถี่ที่ควรทำ
  \item การจัดการผู้ใช้และสิทธิ์ ถ้าระบบมีหน้าจอผู้ดูแล ให้เขียนเป็นขั้นตอนพร้อมภาพแบบคู่มือผู้ใช้
  \item การสำรองและกู้คืนข้อมูล ใส่เฉพาะคำสั่งที่ทดลองกู้คืนได้จริงแล้ว
  \item การปรับรุ่นระบบ การตรวจสถานะ และตำแหน่งบันทึกการทำงาน (log)
  \item ข้อควรระวังด้านความปลอดภัยก่อนและหลังเปิดใช้งานจริง
\end{itemize}
ระบบที่ไม่มีเครื่องแม่ข่ายหรือฐานข้อมูลกลาง เช่น โปรแกรมที่ทำงานบนเครื่องผู้ใช้อย่างเดียว ลบบทนี้ได้
แล้วย้ายเรื่องการปรับรุ่นไปไว้ในบทที่ 2
\end{guide}

\section{งานประจำของผู้ดูแลระบบ}
\label{sec:admin-tasks}

งานที่ผู้ดูแลระบบต้องทำเป็นประจำแสดงดังตารางที่~\ref{tab:admin-tasks}

\begin{table}[H]
\caption{งานประจำของผู้ดูแลระบบ}
\label{tab:admin-tasks}
\centering
\begin{tabular}{|L{\dimexpr\linewidth-7cm-6\tabcolsep-4\arrayrulewidth\relax}|L{3.6cm}|C{3.4cm}|}
\hline
\textbf{งาน} & \textbf{ความถี่} & \textbf{หัวข้อ} \\ \hline
สำรองฐานข้อมูลและไฟล์ของโครงงาน & ทุกวัน & \ref{sec:backup} \\ \hline
ตรวจสถานะระบบและพื้นที่ดิสก์ & ทุกสัปดาห์ & \ref{sec:monitor} \\ \hline
ตรวจบันทึกการทำงานเมื่อผู้ใช้แจ้งปัญหา & เมื่อมีปัญหา & \ref{sec:monitor} \\ \hline
กำหนดบทบาทผู้ดูแลรายวิชาและตั้งรหัสผ่านใหม่ให้ผู้ใช้ & เมื่อมีคำขอ & \ref{sec:users} \\ \hline
ปรับรุ่นระบบ & เมื่อมีรุ่นใหม่ & \ref{sec:upgrade} \\ \hline
\end{tabular}
\end{table}

คำสั่งในบทนี้รันที่โฟลเดอร์ของระบบบนเครื่องแม่ข่าย (\filename{~/cskkuproject}) ตามที่ติดตั้งในบทที่~\ref{ch:install}

\section{การจัดการผู้ใช้}
\label{sec:users}

ผู้ใช้สมัครสมาชิกเองด้วยอีเมลของมหาวิทยาลัย (หัวข้อ~\ref{sec:register})
ระบบมี 3 บทบาท ได้แก่ \cmd{student} (นักศึกษา) \cmd{advisor} (อาจารย์ที่ปรึกษา)
และ \cmd{instructor} (ผู้ดูแลรายวิชา) ระบบรุ่นนี้ยังไม่มีหน้าจอสำหรับผู้ดูแลระบบ
จึงจัดการผู้ใช้ด้วยคำสั่งในหัวข้อย่อยต่อไปนี้

\subsection{การดูรายชื่อผู้ใช้}

\begin{terminal}
docker-compose exec api node scripts/listUsers.js
\end{terminal}

ผลลัพธ์แสดงรหัส อีเมล บทบาท และชื่อของผู้ใช้ทุกคน หนึ่งบรรทัดต่อหนึ่งคน

\subsection{การกำหนดบทบาทผู้ดูแลรายวิชา}

อาจารย์ที่สมัครสมาชิกจะได้บทบาทตามที่เลือกตอนสมัคร
ถ้าต้องการให้อาจารย์เป็นผู้ดูแลรายวิชา ให้เปลี่ยนบทบาทเป็น \cmd{instructor} ด้วยคำสั่งด้านล่าง
(เปลี่ยนอีเมลเป็นของอาจารย์ท่านนั้น)

\begin{terminal}
docker-compose exec db psql -U cskkuadmin cskkuproject
\end{terminal}

จากนั้นพิมพ์คำสั่ง SQL ด้านล่างที่ตัวพร้อมรับคำสั่ง \cmd{cskkuproject=\#} แล้วพิมพ์ \cmd{\q} เพื่อออก

\begin{terminal}
UPDATE users SET role = 'instructor'
  WHERE email = 'teacher@kku.ac.th';
\end{terminal}

\begin{note}
บทบาทใหม่มีผลเมื่อผู้ใช้ออกจากระบบแล้วเข้าสู่ระบบอีกครั้ง เพราะบทบาทถูกบันทึกไว้ใน token ตอนเข้าสู่ระบบ
\end{note}

\subsection{การตั้งรหัสผ่านใหม่ให้ผู้ใช้}

เมื่อผู้ใช้ลืมรหัสผ่าน ผู้ดูแลระบบตั้งรหัสผ่านชั่วคราวให้ได้ด้วยคำสั่งด้านล่าง
ระบบเข้ารหัสรหัสผ่านใหม่ให้อัตโนมัติก่อนบันทึก แล้วแจ้งให้ผู้ใช้เปลี่ยนรหัสผ่านเองที่หน้าโปรไฟล์

\begin{terminal}
docker-compose exec api node -e "
const { User } = require('./models');
User.findOne({ where: { email: 'student@kku.ac.th' } })
  .then((u) => { u.password_hash = 'TempPass2569'; return u.save(); })
  .then(() => process.exit(0));"
\end{terminal}

\section{การสำรองและกู้คืนข้อมูล}
\label{sec:backup}

ข้อมูลของระบบเก็บไว้ 2 แห่ง ได้แก่ ฐานข้อมูล PostgreSQL ซึ่งเก็บผู้ใช้ ข้อเสนอ เนื้อหาไฟล์เอกสาร
ความเห็น และการส่งงาน และ volume \cmd{project_data} ซึ่งเก็บไฟล์ของโครงงานที่ใช้ Compile
ต้องสำรองทั้งสองแห่งทุกครั้ง

\subsection{การสำรองข้อมูล}

สร้างไฟล์สคริปต์ \filename{~/backup.sh} ตามด้านล่าง แล้วให้สิทธิ์รันด้วยคำสั่ง \cmd{chmod +x ~/backup.sh}
ทุกครั้งที่รันสคริปต์ จะได้ไฟล์สำรองที่มีวันที่อยู่ในชื่อไฟล์ในโฟลเดอร์ \filename{~/backup}

\begin{codefile}[language=bash]{~/backup.sh}
#!/bin/sh
mkdir -p ~/backup
cd ~/cskkuproject
# ฐานข้อมูล
docker-compose exec -T db \
  pg_dump -U cskkuadmin -Fc cskkuproject \
  > ~/backup/cskku-$(date +%F).dump
# ไฟล์ของโครงงาน
docker run --rm -v cskkuproject_project_data:/data \
  -v ~/backup:/backup alpine \
  tar czf /backup/files-$(date +%F).tar.gz -C /data .
\end{codefile}

\begin{tip}
ตั้งเวลาสำรองข้อมูลอัตโนมัติทุกคืนเวลา 02:00 น. ด้วยคำสั่ง \cmd{crontab -e} แล้วเพิ่มบรรทัด
\cmd{0 2 * * * ~/backup.sh}
\end{tip}

\subsection{การกู้คืนข้อมูล}

\begin{steps}
  \item หยุดส่วนหลังเพื่อไม่ให้มีผู้ใช้แก้ไขข้อมูลระหว่างกู้คืน
\begin{terminal}
docker-compose stop api
\end{terminal}
  \item กู้คืนฐานข้อมูลจากไฟล์สำรอง (เปลี่ยนชื่อไฟล์ตามวันที่ที่ต้องการ)
\begin{terminal}
docker-compose exec -T db \
  pg_restore -U cskkuadmin -d cskkuproject --clean \
  < ~/backup/cskku-2026-03-01.dump
\end{terminal}
  \item เริ่มส่วนหลังอีกครั้ง แล้วตรวจสอบตามหัวข้อ~\ref{sec:verify}
\begin{terminal}
docker-compose start api
\end{terminal}
\end{steps}

\begin{warning}
การกู้คืนเขียนทับข้อมูลปัจจุบันทั้งหมด ทดลองกู้คืนบนเครื่องทดสอบก่อนเสมอ
และสุ่มทดลองกู้คืนไฟล์สำรองเป็นระยะเพื่อยืนยันว่าไฟล์ใช้งานได้
\end{warning}

\section{การปรับรุ่นระบบ}
\label{sec:upgrade}

\begin{steps}
  \item สำรองข้อมูลตามหัวข้อ~\ref{sec:backup}
  \item ดาวน์โหลดซอร์สโค้ดรุ่นใหม่
\begin{terminal}
git pull
\end{terminal}
  \item สร้างและเริ่มส่วนหลังรุ่นใหม่
\begin{terminal}
docker-compose up -d --build api
\end{terminal}
  \item Build ส่วนหน้าใหม่และคัดลอกไปยังโฟลเดอร์ของเว็บเซิร์ฟเวอร์ ตามขั้นตอนที่ 6 ในหัวข้อ~\ref{sec:install-steps}
  \item ตรวจสอบระบบตามหัวข้อ~\ref{sec:verify}
\end{steps}

\begin{important}
ส่วนหลังสร้างตารางที่ยังไม่มีในฐานข้อมูลให้อัตโนมัติเมื่อเริ่มทำงาน แต่ไม่แก้ไขตารางเดิม
ถ้ารุ่นใหม่เปลี่ยนโครงสร้างตาราง ให้ทำตามบันทึกการเปลี่ยนแปลงของรุ่นนั้นก่อนเริ่มส่วนหลัง
\end{important}

\section{การตรวจสถานะและบันทึกการทำงาน}
\label{sec:monitor}

ตรวจสถานะของระบบด้วยคำสั่งในตารางที่~\ref{tab:monitor}
ส่วนหลังบันทึกคำขอทุกรายการและข้อผิดพลาดไว้ในบันทึกของคอนเทนเนอร์ \cmd{api}

\begin{table}[H]
\caption{การตรวจสถานะของระบบ}
\label{tab:monitor}
\centering
\begin{tabular}{|L{3.4cm}|L{\dimexpr\linewidth-7.6cm-6\tabcolsep-4\arrayrulewidth\relax}|L{4.2cm}|}
\hline
\textbf{สิ่งที่ตรวจ} & \textbf{คำสั่ง} & \textbf{ผลที่ปกติ} \\ \hline
สถานะบริการ & \cmd{docker-compose ps} & \cmd{db} และ \cmd{api} มีสถานะ \cmd{Up} \\ \hline
การตอบสนองของส่วนหลัง & \cmd{curl http://localhost:3001/api/health} & \cmd{\{"status":"ok"} \\ \hline
บันทึกการทำงาน & \cmd{docker-compose logs --tail=100 api} & ไม่มีข้อความ \cmd{Error} ซ้ำ ๆ \\ \hline
พื้นที่ดิสก์ & \cmd{df -h} และ \cmd{docker system df} & พื้นที่ว่างมากกว่าร้อยละ 20 \\ \hline
\end{tabular}
\end{table}

\section{ข้อควรระวังด้านความปลอดภัย}
\label{sec:security}

\begin{enumerate}
  \item ตั้งค่า \cmd{JWT_SECRET} และ \cmd{DB_PASSWORD} ใหม่ ห้ามใช้ค่าตัวอย่าง
  \item ตั้ง \cmd{NODE_ENV=production} ระบบจะไม่แสดงรายละเอียดข้อผิดพลาดภายในให้ผู้ใช้เห็น
  \item ลบหรือเปลี่ยนรหัสผ่านของบัญชีทดสอบ (ตารางที่~\ref{tab:demo-accounts})
  \item เปิด HTTPS ที่ Nginx เช่น ใช้ใบรับรองจาก Let's Encrypt
  \item ปิดพอร์ต 3001 และ 5432 จากภายนอกด้วยไฟร์วอลล์ ให้เข้าระบบผ่าน Nginx เท่านั้น
  \item คอนเทนเนอร์ที่ใช้ Compile ปิด shell escape และไม่มีการเชื่อมต่อเครือข่าย ห้ามเปลี่ยนค่าเหล่านี้
\end{enumerate}
